% !TITLE 静夜思
% !AUTHOR 李白
% !DYNASTY 唐
% !ORDER 6

\begin{poem}
床前明月光，疑\textsuperscript{1}是地上霜。
举头望明月，低头思故乡。
\end{poem}

\begin{annotations}
\notetext{1}{疑：怀疑、恍惚以为是。月光照在地上，恍惚间以为是白霜。“疑”字写出人在半梦半醒之间感官的模糊。}
\end{annotations}

\begin{appreciation}
《静夜思》是中国流传最广的诗之一。二十个字，没有生僻字，没有典故，每个中国人都能背。可就是这样一首诗，一千多年了，还不断有人写文章讨论：它到底好在哪里？

先看那个“疑”字：“床前明月光，疑是地上霜。”疑是怀疑、恍惚。诗人不是看见月光，而是恍惚间以为是白霜。为什么是霜不是雪？霜是秋天的，凉的，天亮就化的。一个独处的深夜，四周清冷，人半睡半醒，看到一地白光，第一个念头就是“是不是下霜了”——这是被凉意叫醒的人才会有的错觉。一个“疑”字，写尽了深夜的寂静和人的恍惚。

“举头望明月，低头思故乡。”两个动作，一仰一俯，没有任何解释。全诗没有一个形容词，全是动作：疑、举头、望、低头、思。李白不告诉你他为什么想家，也不说想家的滋味，他只给你看这两个动作。可就是这两个动作，让每一个离家的人心里一动——因为想家的人，都做过一模一样的动作。

还记得《行行重行行》吗？两千年前有个女子，想念离家的游子，想得“衣带日已缓”——衣带一天比一天松，人瘦了。她盼着游子回头，最后只说了一句“努力加餐饭”。到了李白，思念换了个办法：不用瘦，不用叮嘱，抬头看一眼月亮就够了。两千年过去，想家的办法没变过——不是看月亮，就是好好吃饭。

你以后去了外地，想家的时候抬头看看月亮就行。月亮照得到你，也就照得到家。李白当年也是这么看的。
\end{appreciation}
